\section{Closure deficit and a variational obstruction}\label{sec:closure}

This section proves Theorem~\ref{thm:NG_MACRO_CLOSURE_DEFICIT}.
The setting is finite throughout.
Let $X$ be a finite state space, let $Y$ be a finite macro space, and let $(X_t)_{t\in\mathbb{Z}}$
be a stationary Markov chain on $X$ with stationary law $\mu$.
Fix a deterministic packaging map $\Pi:X\to Y$ and define the macro process by
$Y_t := \Pi(X_t)$.
For an integer lag $\tau\ge 1$ and a microstate $x\in X$, write
\[
p_x^{(\tau)}(\cdot) := \mathrm{Law}(Y_{t+\tau}\mid X_t=x)\in \Delta(Y).
\]
For a macrostate $y\in Y$, write
\[
\bar p_y^{(\tau)}(\cdot) := \mathrm{Law}(Y_{t+\tau}\mid Y_t=y)\in \Delta(Y).
\]
A macro kernel is a function $K:Y\to \Delta(Y)$, equivalently a row-stochastic array
$K(y,y')$ on $Y\times Y$.
The variational objective at lag $\tau$ is
\[
\mathcal L_\tau(K) := \sum_{x\in X} \mu(x)\,D_{\mathrm{KL}}\!\left(p_x^{(\tau)}\,\middle\|\,K(\Pi(x),\cdot)\right).
\]
The KL decompositions and conditional-mutual-information identities used below are standard finite information-theoretic calculations, and the variational minimizer is a finite KL-projection instance \citep{CoverThomas2006,Csiszar1975}.

\begin{lemma}[KL mixture decomposition]\label{lem:closure_kl_mixture}
Let $(w_i)_{i\in I}$ be nonnegative weights on a finite index set $I$, not all zero,
and let $(P_i)_{i\in I}$ be probability laws on a finite set $\Omega$.
Define
\[
W := \sum_{i\in I} w_i,
\qquad
R := \frac{1}{W}\sum_{i\in I} w_i P_i.
\]
Then for every probability law $Q$ on $\Omega$,
\[
\sum_{i\in I} w_i\,D_{\mathrm{KL}}\!\left(P_i\,\middle\|\,Q\right)
=
\sum_{i\in I} w_i\,D_{\mathrm{KL}}\!\left(P_i\,\middle\|\,R\right)
+
W\,D_{\mathrm{KL}}\!\left(R\,\middle\|\,Q\right).
\]
In particular, $Q=R$ is a minimizer, and the minimum value is
\[
\sum_{i\in I} w_i\,D_{\mathrm{KL}}\!\left(P_i\,\middle\|\,R\right).
\]
\end{lemma}

\begin{proof}
See Appendix~\ref{app:varclosure}.
\end{proof}

\begin{lemma}[Best macro kernel]\label{lem:closure_best_macro_kernel}
For each $y\in Y$, define
\[
K^\star_\tau(y,\cdot) := \bar p_y^{(\tau)}(\cdot)=\mathrm{Law}(Y_{t+\tau}\mid Y_t=y).
\]
Then
\[
\min_K \mathcal L_\tau(K)
=
\sum_{x\in X}\mu(x)\,D_{\mathrm{KL}}\!\left(p_x^{(\tau)}\,\middle\|\,\bar p_{\Pi(x)}^{(\tau)}\right),
\]
and the minimum is attained by $K^\star_\tau$.
\end{lemma}

\begin{proof}
See Appendix~\ref{app:varclosure}.
\end{proof}

\begin{lemma}[Conditional mutual information identity]\label{lem:closure_cmi_identity}
In this finite deterministic setting,
\[
I(X_t;Y_{t+\tau}\mid Y_t)
=
\sum_{x\in X}\mu(x)\,D_{\mathrm{KL}}\!\left(p_x^{(\tau)}\,\middle\|\,\bar p_{\Pi(x)}^{(\tau)}\right).
\]
\end{lemma}

\begin{proof}
See Appendix~\ref{app:varclosure}.
\end{proof}

\begin{proof}[Proof of Theorem~\ref{thm:NG_MACRO_CLOSURE_DEFICIT}]
By definition, the closure deficit in the theorem is
\[
\mathrm{CD}_\tau(\Pi)
=
I(X_t;\Pi(X_{t+\tau})\mid \Pi(X_t))
=
I(X_t;Y_{t+\tau}\mid Y_t).
\]
Lemma~\ref{lem:closure_cmi_identity} expresses this quantity as
\[
\mathrm{CD}_\tau(\Pi)
=
\sum_{x\in X}\mu(x)\,D_{\mathrm{KL}}\!\left(p_x^{(\tau)}\,\middle\|\,\bar p_{\Pi(x)}^{(\tau)}\right).
\]
Lemma~\ref{lem:closure_best_macro_kernel} identifies the same expression as the minimum
of the variational objective:
\[
\mathrm{CD}_\tau(\Pi)
=
\min_K \mathcal L_\tau(K),
\qquad
K^\star_\tau(y,\cdot)=\bar p_y^{(\tau)}(\cdot).
\]
This proves the finite variational identity and the formula for the minimizing macro kernel.

For the positivity claim, suppose there exist $x,x'\in X$ with $\Pi(x)=\Pi(x')=:y$,
$\mu(x)>0$, $\mu(x')>0$, and $p_x^{(\tau)}\neq p_{x'}^{(\tau)}$.
Then not all positive-mass fiber laws in $\Pi^{-1}(y)$ coincide with their weighted average
$\bar p_y^{(\tau)}$.
Hence at least one state $z\in\Pi^{-1}(y)$ with $\mu(z)>0$ satisfies
\[
D_{\mathrm{KL}}\!\left(p_z^{(\tau)}\,\middle\|\,\bar p_y^{(\tau)}\right)>0.
\]
The sum formula for $\mathrm{CD}_\tau(\Pi)$ therefore contains at least one strictly
positive term and all remaining terms are nonnegative, so
\[
\mathrm{CD}_\tau(\Pi)>0.
\]
\end{proof}

\begin{remark}[Scope of exact closure]
If every macro fiber is a singleton, then $\Pi$ loses no information and the closure
deficit vanishes trivially.
More generally, if all packaged future laws $p_x^{(\tau)}$ agree inside each fiber
$\Pi^{-1}(y)$, then the minimizing macro kernel reproduces those common laws and
$\mathrm{CD}_\tau(\Pi)=0$.
The theorem becomes informative exactly when a macro fiber contains distinct packaged
future laws.
\end{remark}

\begin{remark}[Attribution is separate]
The theorem isolates a finite obstruction: distinct packaged future laws inside a single
macro fiber force a positive closure deficit.
Which primitive or modeling decision produces such a pair is a separate question and is
independent of the variational identity proved here.
\end{remark}
